\documentclass[11pt,a4paper]{article}
\usepackage[margin=1in]{geometry}
\usepackage{amsmath,amssymb}
\usepackage{booktabs}
\usepackage{enumitem}
\usepackage{hyperref}
\usepackage{xcolor}
\usepackage{graphicx}
\usepackage{multirow}
\usepackage{array}
\usepackage{longtable}

\definecolor{passing}{HTML}{2E7D32}
\definecolor{failing}{HTML}{C62828}
\definecolor{partial}{HTML}{EF6C00}

\title{Achieving RGB Depth-Estimation Student Convergence\\on Gap Terrain: Experiment Plan}
\author{Parkour Student Distillation Project}
\date{March 2026}

\begin{document}
\maketitle

%----------------------------------------------------------------------
\section{Problem Statement}
%----------------------------------------------------------------------

A Go2 quadruped is trained to traverse parkour terrain (stairs, hurdles,
gaps) via DAgger-based teacher--student distillation.  The \emph{teacher}
uses ground-truth (GT) scandot heightmaps and achieves 100\% success on all
terrain families.  Two student variants receive only front-facing depth
images through a shared CNN$\,+\,$GRU encoder:

\begin{enumerate}[nosep]
  \item \textbf{GT depth student} --- uses simulator-rendered depth.
        Currently at \textbf{100\% success on all families} (step 2939/5000).
  \item \textbf{RGB depth-estimation student} --- uses Depth Anything V2
        (DA-V2) inferred from RGB.  Currently at \textbf{0\% gap success}
        (step 1088), despite 100\% stairs and 79\% hurdles.
\end{enumerate}

The gap between the two students must be closed.  This document enumerates
every hypothesis, proposes targeted experiments, and defines a rigorous
methodology for parallel execution across 4 GPUs.

%----------------------------------------------------------------------
\section{Background and Current Architecture}
%----------------------------------------------------------------------

\subsection{DAgger Distillation Pipeline}

\begin{itemize}[nosep]
  \item \textbf{Teacher}: PPO-trained, proprioception (53) $+$ scandots (132)
        $\to$ MLP [512,\,256,\,128] $\to$ 12 joint actions.
  \item \textbf{Student depth encoder}: depth image $(1, 58, 87)$ $\to$
        CNN $\to$ 32-dim $+$ masked proprio (53) $\to$ MLP [85,\,128,\,32]
        $\to$ GRU (512 hidden) $\to$ MLP $\to$ latent (32) $+$ yaw (2).
  \item \textbf{Student actor}: proprio (53, with predicted yaw injected via
        MTS) $+$ latent (32) $\to$ MLP [512,\,256,\,128] $\to$ 12 actions.
  \item \textbf{Loss}: $\mathcal{L} = \|a_{\text{teacher}} - a_{\text{student}}\|_2
        + \|y_{\text{oracle}} - y_{\text{pred}}\|_2$, windowed BPTT
        (window $= 24$ steps) through GRU, Adam LR $= 2 \times 10^{-3}$.
\end{itemize}

\subsection{Depth Normalization Pipelines}

\textbf{GT depth} (simulator $\to$ student):
\[
  d_{\text{norm}} = \frac{\text{clamp}(d_{\text{raw}},\; 0,\; 2)}{2} - 0.5
  \quad\Rightarrow\quad d_{\text{norm}} \in [-0.5,\; 0.5]
\]
Fixed, absolute mapping.  Mean $\approx 0.41$ (most pixels at far clip).

\textbf{Inferred depth} (DA-V2 metric outdoor $\to$ EMA normalization):
\[
  d_{\text{norm}} = \frac{d_{\text{raw}} - \text{EMA}_{p2}}
                         {\text{EMA}_{p98} - \text{EMA}_{p2}} - 0.5
  \quad\Rightarrow\quad d_{\text{norm}} \in [-0.5,\; 0.5]
\]
Adaptive, EMA $\alpha = 0.02$ over batch percentiles.  Mean $\approx 0.0$.

\subsection{Measured Depth Estimation Quality}

\begin{table}[h]
\centering
\begin{tabular}{lcccc}
\toprule
\textbf{Model} & \textbf{Raw Range} & \textbf{Spatial $r$} & \textbf{Direction} & \textbf{Affine MAE} \\
\midrule
GT depth (simulator)            & 0.45--2.0\,m & 1.00 & --- & 0.000 \\
DA-V2 Relative Small            & 0.4--5.4 (disp.)    & $-0.59$ & Inverted & 0.110 \\
DA-V2 Metric Indoor Small       & 1.0--4.5\,m  & $+0.72$ & Correct & 0.093 \\
DA-V2 Metric Outdoor Small      & 4.6--21.0\,m & $+0.79$ & Correct & 0.091 \\
\bottomrule
\end{tabular}
\caption{Depth estimation model comparison on Genesis synthetic renders (106$\times$60).
$r$ = Pearson correlation with GT depth per-pixel. Affine MAE after optimal least-squares alignment.}
\label{tab:depth_models}
\end{table}

Key finding: ``Metric'' models are \emph{not} metric-accurate on synthetic
renders --- the outdoor model predicts 5--10$\times$ too far.  However, it
has the best \emph{spatial correlation} ($r = 0.79$), meaning it preserves
relative depth structure.

%----------------------------------------------------------------------
\section{Hypotheses}
%----------------------------------------------------------------------

We identify seven hypotheses for why the RGB student fails on gaps, ordered
by estimated likelihood and ease of testing:

\begin{enumerate}[label=\textbf{H\arabic*}]
  \item \textbf{Normalization distribution mismatch.}
        EMA normalization maps to mean $\approx 0$ with uniform spread, while
        GT maps to mean $\approx 0.41$ with heavy right skew (most pixels at
        far clip).  The CNN+GRU may need the skewed distribution to reliably
        detect gaps as ``outlier low'' pixels.

  \item \textbf{Depth estimation spatial resolution insufficient for gaps.}
        DA-V2 at 106$\times$60 input may not resolve the narrow gap geometry
        ($\sim$5--10 pixel wide stripe at 0.5\,m distance).  Larger models
        (base/large) or higher resolution may help.

  \item \textbf{Depth estimation update frequency too low.}
        Current \texttt{update\_interval=2} means the depth model runs every
        2nd sensor step (every 10 physics steps = 200\,ms).  Gap crossing
        requires precise timing; stale depth may cause missed decisions.

  \item \textbf{BPTT window too short for noisy depth.}
        With noisier depth, the GRU needs more temporal context to filter
        signal from noise.  Window = 24 may be insufficient; extreme-parkour
        uses similar but with cleaner GT depth.

  \item \textbf{Learning rate / optimizer dynamics.}
        The depth encoder introduces more parameters with noisier gradients.
        A lower LR or warmup schedule may stabilize training.

  \item \textbf{Depth noise level too high.}
        \texttt{depth\_noise\_level = 0.1} adds uniform noise $\pm 0.1$ to
        normalized depth.  Combined with estimation noise, total noise may
        overwhelm the gap signal (gaps are $\sim$0.1--0.2 normalized units
        different from background).

  \item \textbf{Depth model wrong entirely.}
        Perhaps a different backbone (Metric3D, UniDepth, Video Depth
        Anything) would give better spatial structure on synthetic renders.
\end{enumerate}

%----------------------------------------------------------------------
\section{Experimental Design}
%----------------------------------------------------------------------

\subsection{Phase 0: Oracle Experiments (Critical Controls)}

These experiments isolate the contribution of each component.  They must
run \emph{first} because they determine which subsequent experiments are
worth running.

\begin{table}[h]
\centering
\small
\begin{tabular}{p{2.5cm}p{4cm}p{5.5cm}p{2.5cm}}
\toprule
\textbf{ID} & \textbf{Modification} & \textbf{What it tests} & \textbf{Baseline} \\
\midrule
\texttt{O1-gt-ema} &
  GT depth $+$ EMA normalization (bypass simulator linear norm, apply EMA to raw GT meters) &
  Is the EMA normalization itself the problem, independent of the depth model? If GT+EMA also fails on gaps, normalization is the bottleneck. &
  GT student \\
\midrule
\texttt{O2-gt-affine} &
  GT depth $+$ random affine perturbation: $d' = (1 + \epsilon_s)\,d + \epsilon_b$, $\epsilon_s \sim \mathcal{U}(-0.3, 0.3)$, $\epsilon_b \sim \mathcal{U}(-0.2, 0.2)$ per frame &
  How robust is the CNN to distribution shift? If GT+affine still works, the CNN can tolerate distribution differences. &
  GT student \\
\midrule
\texttt{O3-rgb-gtfallback} &
  RGB student with \texttt{update\_interval=1} (depth estimation every sensor step) &
  Is the 2$\times$ update interval causing stale depth on gaps? &
  Current RGB \\
\bottomrule
\end{tabular}
\caption{Phase 0 oracle experiments.}
\label{tab:phase0}
\end{table}

\subsection{Phase 1: Normalization Sweep}

Conditioned on Phase 0 results.  If \texttt{O1-gt-ema} fails, normalization
is the bottleneck and these are high priority.

\begin{table}[h]
\centering
\small
\begin{tabular}{p{2.5cm}p{5cm}p{5cm}}
\toprule
\textbf{ID} & \textbf{Configuration} & \textbf{Rationale} \\
\midrule
\texttt{N1-affine-fit} &
  Pre-computed global affine: $d_{\text{norm}} = 0.088 \cdot d_{\text{raw}} + 0.429 - 0.5$,
  then clamp to $[-0.5, 0.5]$ (from least-squares fit) &
  Fixed mapping avoids EMA adaptation lag.  Uses the measured optimal affine
  alignment to GT. \\
\midrule
\texttt{N2-ema-fast} &
  EMA with $\alpha = 0.1$ (5$\times$ faster adaptation) &
  Current $\alpha = 0.02$ may be too slow to track scene changes during
  curriculum progression. \\
\midrule
\texttt{N3-ema-slow} &
  EMA with $\alpha = 0.005$ (4$\times$ slower) &
  Slower EMA = more stable normalization; tests if EMA drift is the issue. \\
\midrule
\texttt{N4-quantile-wide} &
  EMA on $p_1 / p_{99}$ instead of $p_2 / p_{98}$ &
  Wider percentile bounds preserve more of the distribution tails, which
  may contain the gap signal. \\
\midrule
\texttt{N5-running-meanstd} &
  Normalize as $(d - \mu_{\text{EMA}}) / \sigma_{\text{EMA}}$, then
  $\tanh / 2$ to map to $[-0.5, 0.5]$ &
  Mean/std normalization (like BatchNorm) instead of percentile; may better
  preserve gap outlier structure. \\
\bottomrule
\end{tabular}
\caption{Phase 1 normalization experiments.}
\label{tab:phase1}
\end{table}

\subsection{Phase 2: Depth Model Sweep}

If Phase 0 shows the depth model is the bottleneck (i.e., \texttt{O1-gt-ema}
succeeds), these experiments test alternative depth estimation backends.

\begin{table}[h]
\centering
\small
\begin{tabular}{p{2.5cm}p{5cm}p{5cm}}
\toprule
\textbf{ID} & \textbf{Model} & \textbf{Rationale} \\
\midrule
\texttt{D1-indoor} &
  DA-V2 Metric Indoor Small ($r = 0.72$, range 1.0--4.5\,m) &
  Closer absolute range to GT than outdoor; may work better with fixed
  affine normalization. \\
\midrule
\texttt{D2-relative-inv} &
  DA-V2 Relative Small with inversion ($d' = d_{\max} - d$) + EMA norm &
  Relative model has inverted output; after inversion, $r \approx +0.59$.
  Tests if ordinal depth with correct direction suffices. \\
\midrule
\texttt{D3-outdoor-base} &
  DA-V2 Metric Outdoor \textbf{Base} ($r$ likely $> 0.79$) &
  Larger encoder may give better spatial resolution for gap detection.
  Slower but more accurate. \\
\midrule
\texttt{D4-video-metric} &
  Video Depth Anything Metric Small (temporal consistency built-in) &
  Uses temporal priors across frames; may give smoother depth for the GRU.
  Note: streaming per-env state, higher latency. \\
\bottomrule
\end{tabular}
\caption{Phase 2 depth model experiments.}
\label{tab:phase2}
\end{table}

\subsection{Phase 3: Training Hyperparameter Sweep}

If the depth signal is adequate (Phases 0--2 confirm), these experiments
tune the learning dynamics.

\begin{table}[h]
\centering
\small
\begin{tabular}{p{2.5cm}p{5cm}p{5cm}}
\toprule
\textbf{ID} & \textbf{Configuration} & \textbf{Rationale} \\
\midrule
\texttt{T1-bptt32} &
  BPTT window = 32 (from 24) &
  Longer temporal gradient flow through GRU; helps if the depth encoder
  needs more context to learn gap detection from noisy input. \\
\midrule
\texttt{T2-bptt48} &
  BPTT window = 48 &
  Even longer window; tests diminishing returns vs.\ memory cost. \\
\midrule
\texttt{T3-lr-low} &
  Learning rate $= 5 \times 10^{-4}$ (4$\times$ lower) &
  Slower, more stable optimization; may prevent the oscillation pattern
  where gap success spikes then collapses. \\
\midrule
\texttt{T4-lr-warmup} &
  LR warmup: linear 0 $\to$ $2 \times 10^{-3}$ over first 200 iterations &
  Prevents early instability while EMA normalization is still calibrating.\\
\midrule
\texttt{T5-noise-zero} &
  \texttt{depth\_noise\_level = 0.0} &
  Removes additive depth noise entirely.  Tests if 0.1 noise overwhelms
  the already-noisy inferred depth signal. \\
\midrule
\texttt{T6-noise-low} &
  \texttt{depth\_noise\_level = 0.02} &
  Reduced noise: 5$\times$ lower than current 0.1. \\
\midrule
\texttt{T7-steps240} &
  \texttt{num\_steps\_per\_env = 240} (from 120) &
  Longer rollouts give more data per gradient step; may help with
  sparse gap encounters (gaps are $1/3$ of terrain). \\
\midrule
\texttt{T8-update-1} &
  \texttt{update\_interval = 1} (depth est every sensor step) &
  Same as \texttt{O3} but in combination with best normalization
  from Phase 1. \\
\bottomrule
\end{tabular}
\caption{Phase 3 training hyperparameter experiments.}
\label{tab:phase3}
\end{table}

\subsection{Phase 4: Architecture Modifications}

If training dynamics are not the issue, these experiments modify the
student architecture.

\begin{table}[h]
\centering
\small
\begin{tabular}{p{2.5cm}p{5cm}p{5cm}}
\toprule
\textbf{ID} & \textbf{Modification} & \textbf{Rationale} \\
\midrule
\texttt{A1-proprio-hist} &
  Add 10-frame proprioceptive history encoder (1D conv as in
  extreme-parkour) as additional input to actor &
  Provides a non-vision temporal signal.  Gap crossing requires precise
  timing; joint feedback history may compensate for noisy depth. \\
\midrule
\texttt{A2-gru1024} &
  GRU hidden dim = 1024 (from 512) &
  Larger temporal memory.  The GRU must encode both depth features and
  temporal state; 512 may be insufficient with noisy input. \\
\midrule
\texttt{A3-2frame-input} &
  Pass both current \emph{and} previous depth frame as 2-channel input to
  CNN: shape $(2, 58, 87)$ &
  Explicit temporal stacking at CNN level, bypassing GRU for motion cues.
  Requires modifying \texttt{DepthBackbone58x87} input channels. \\
\bottomrule
\end{tabular}
\caption{Phase 4 architecture experiments.}
\label{tab:phase4}
\end{table}

\subsection{Phase 5: Combined Best Configuration}

After Phases 0--4 identify which individual changes help, run a final
experiment combining the top 2--3 improvements.

%----------------------------------------------------------------------
\section{Possible Bugfixes to Investigate}
%----------------------------------------------------------------------

Before running sweeps, we should audit the following potential bugs:

\begin{enumerate}[label=\textbf{B\arabic*}]
  \item \textbf{EMA not reset on curriculum change.}
        When the terrain curriculum advances (harder terrains), the depth
        distribution shifts.  The EMA may lag behind, giving incorrect
        normalization for the new terrain.  \emph{Fix}: Reset EMA state when
        terrain level changes, or increase $\alpha$ during level transitions.

  \item \textbf{Crop mismatch between GT and inferred paths.}
        GT depth is cropped by the \emph{simulator} before normalization.
        Inferred depth is cropped in \texttt{\_process\_inferred\_depth}.
        Verify the crop parameters match exactly and the resulting FOVs are
        identical.  Currently both use
        \texttt{crop\_top=0, bottom=2, left=4, right=4} from
        \texttt{depth\_camera\_config}, which should be correct.

  \item \textbf{Inferred depth buffer timing.}
        The GT path writes to the depth buffer via the simulator's
        \texttt{\_update\_depth\_images} (shift right, insert at 0).
        The inferred path does its own shift in
        \texttt{\_refresh\_student\_depth}.  Verify the timing is identical:
        both should shift \emph{after} the new frame is available, and the
        student should see the \emph{previous} frame (index 1).

  \item \textbf{EMA initialization from first batch.}
        The first batch may be atypical (all envs start from the same initial
        position, seeing similar depth).  This seeds the EMA with a biased
        estimate.  \emph{Fix}: Initialize from the mean of the first $K$
        batches instead of just the first.

  \item \textbf{GRU hidden not reset on env reset for inferred path.}
        Verify that \texttt{reset\_gru(env\_ids)} is called for reset envs
        in the inferred depth path.  If GRU hidden carries stale temporal
        state across episode boundaries, it could confuse gap detection.
        (This should be handled by the runner, but worth auditing.)

  \item \textbf{Depth estimation model receives wrong color format.}
        Genesis may return RGB or BGR; DA-V2 expects RGB.  Verify the color
        channel ordering in \texttt{\_rgb\_batch\_to\_numpy}.
\end{enumerate}

%----------------------------------------------------------------------
\section{Methodology}
%----------------------------------------------------------------------

\subsection{Evaluation Metric}

The primary metric is \textbf{\texttt{Episode/success\_gap}} averaged over
the last 200 training iterations.  Secondary metrics:

\begin{itemize}[nosep]
  \item \texttt{Episode/success\_stairs} --- must not regress below 90\%.
  \item \texttt{Episode/success\_hurdle\_block} --- must not regress below 60\%.
  \item \texttt{Loss/action} --- lower is better; indicates distillation quality.
  \item \texttt{Episode/progress\_gap} --- gap progress ratio, useful when
        success is 0\% (partial credit).
\end{itemize}

\subsection{Success Criteria}

An experiment is considered \textbf{successful} if it achieves:
\begin{itemize}[nosep]
  \item $\texttt{success\_gap} \geq 50\%$ sustained for $\geq 200$ iterations, OR
  \item $\texttt{success\_gap} \geq 30\%$ with a clear upward trend at step 2000.
\end{itemize}

The ultimate target is $\texttt{success\_gap} \geq 80\%$, matching the GT
student's trajectory.

\subsection{Training Protocol}

\begin{itemize}[nosep]
  \item Each experiment runs for \textbf{2000 iterations} (sufficient to see
        gap learning based on GT student trajectory).
  \item \texttt{num\_envs = 48}, \texttt{num\_steps\_per\_env = 120}
        (unless varied by the experiment).
  \item Terrain curriculum enabled (default), all families mixed.
  \item Seed = 1 (default).  If a promising result is found, confirm with
        seed = 42.
  \item Training launched via:\\
        \texttt{CUDA\_VISIBLE\_DEVICES=\$GPU .venv/bin/python
        legged\_gym/scripts/train.py --task go2\_parkour\_depth\_est\_student
        --headless --max\_iterations 2000}
  \item Config overrides applied programmatically before launch (modify
        config class attributes in the training script).
\end{itemize}

\subsection{GPU Allocation}

Four GPUs are available.  Experiments are scheduled in phases with dependencies:

\begin{table}[h]
\centering
\begin{tabular}{ccccc}
\toprule
\textbf{Phase} & \textbf{GPU 0} & \textbf{GPU 1} & \textbf{GPU 2} & \textbf{GPU 3} \\
\midrule
0 & \texttt{O1-gt-ema} & \texttt{O2-gt-affine} & \texttt{O3-rgb-gtfallback} & (continue current run) \\
\midrule
\multicolumn{5}{c}{\emph{Wait for Phase 0 results to determine priority of Phase 1 vs 2}} \\
\midrule
1 & \texttt{N1-affine-fit} & \texttt{N2-ema-fast} & \texttt{N5-running-meanstd} & \texttt{T5-noise-zero} \\
1b & \texttt{N3-ema-slow} & \texttt{N4-quantile-wide} & \texttt{T1-bptt32} & \texttt{T3-lr-low} \\
\midrule
2 & \texttt{D1-indoor} & \texttt{D2-relative-inv} & \texttt{D3-outdoor-base} & \texttt{D4-video-metric} \\
\midrule
3 & \texttt{T2-bptt48} & \texttt{T4-lr-warmup} & \texttt{T6-noise-low} & \texttt{T8-update-1} \\
\midrule
4 & \texttt{A1-proprio-hist} & \texttt{A2-gru1024} & \texttt{A3-2frame-input} & (best combo) \\
\bottomrule
\end{tabular}
\caption{GPU scheduling plan.  Each cell runs for $\sim$2000 iterations.
Phase 0 takes priority; subsequent phases are re-ordered based on results.}
\label{tab:gpu_schedule}
\end{table}

\subsection{Implementation Notes for Config Overrides}

Each experiment requires modifying the config before training.  The cleanest
approach is a wrapper script that:

\begin{enumerate}[nosep]
  \item Imports the config classes.
  \item Applies experiment-specific overrides (e.g., changing EMA alpha,
        model type, BPTT window).
  \item Calls the standard training entry point.
\end{enumerate}

For overrides that require code changes (e.g., affine normalization,
mean/std normalization, proprio history encoder), a git worktree per
experiment keeps changes isolated.

\subsection{Monitoring and Early Stopping}

A monitoring script reads TensorBoard events every 5 minutes and reports:
\begin{itemize}[nosep]
  \item Current step, \texttt{success\_gap}, \texttt{progress\_gap},
        \texttt{action\_loss} for each running experiment.
  \item Early stop if \texttt{success\_stairs} drops below 50\% for 100
        consecutive iterations (training has diverged).
  \item Highlight any experiment where \texttt{success\_gap} $> 20\%$
        (promising signal).
\end{itemize}

%----------------------------------------------------------------------
\section{Execution Priority}
%----------------------------------------------------------------------

Based on the evidence so far, the recommended execution order is:

\begin{enumerate}
  \item \textbf{Phase 0 (oracle)} --- \emph{Must run first.}
        \texttt{O1-gt-ema} is the single most informative experiment.
        If GT depth + EMA normalization works, the problem is purely depth
        model quality.  If it fails, the problem is normalization.

  \item \textbf{Bugfix audit (B1--B6)} --- Quick code review, no GPU needed.
        Do in parallel with Phase 0.

  \item \textbf{Phase 1 or 2} --- Conditioned on Phase 0:
    \begin{itemize}[nosep]
      \item If \texttt{O1-gt-ema} \textbf{fails}: normalization is the
            bottleneck $\to$ prioritize Phase 1.
      \item If \texttt{O1-gt-ema} \textbf{succeeds}: depth model is the
            bottleneck $\to$ prioritize Phase 2.
    \end{itemize}

  \item \textbf{Phase 3} --- Training hyperparameters.  Run the most
        promising 4 in parallel after identifying the best normalization/model.

  \item \textbf{Phase 4} --- Architecture changes.  Only if Phases 1--3
        are insufficient.

  \item \textbf{Phase 5} --- Combine the top 2--3 improvements from all
        phases into a single run.
\end{enumerate}

%----------------------------------------------------------------------
\section{Risk Assessment}
%----------------------------------------------------------------------

\begin{itemize}
  \item \textbf{High risk}: The fundamental limit may be DA-V2's spatial
        resolution on 106$\times$60 synthetic images.  Gap features are
        $\sim$5--10 pixels wide at decision distance.  If no depth model
        resolves gaps, we may need to increase camera resolution (but this
        slows training).

  \item \textbf{Medium risk}: The oscillation pattern (gap success spikes
        then collapses) suggests a training stability issue, not a
        fundamental information deficit.  This is encouraging --- the
        student \emph{briefly} learns gaps, then forgets.  LR scheduling or
        longer BPTT may suffice.

  \item \textbf{Low risk}: The curriculum may systematically under-sample
        gap terrain, giving the optimizer too few gap gradients.  Could
        address with curriculum weighting, but this is unlikely given the
        teacher converges fine.
\end{itemize}

%----------------------------------------------------------------------
\section{Expected Timeline}
%----------------------------------------------------------------------

\begin{itemize}[nosep]
  \item \textbf{Phase 0}: 2000 iterations $\approx$ 2--3 hours per run,
        3 runs in parallel.
  \item \textbf{Phase 1--2}: 4 runs in parallel, 2--3 hours each; two
        sub-rounds = 4--6 hours.
  \item \textbf{Phase 3}: 4 runs in parallel, 2--3 hours.
  \item \textbf{Phase 4}: 3 runs, 2--3 hours.
  \item \textbf{Phase 5}: 1 final run, 5000 iterations $\approx$ 5--6 hours.
  \item \textbf{Total}: $\sim$12--18 hours of wall time with 4 GPUs.
\end{itemize}

%----------------------------------------------------------------------
\section{Appendix: Experiment Implementation Details}
%----------------------------------------------------------------------

\subsection{\texttt{O1-gt-ema}: GT Depth with EMA Normalization}

Override the GT student config to use EMA normalization instead of
simulator linear normalization.  This requires:
\begin{enumerate}[nosep]
  \item Set \texttt{add\_depth = True}, \texttt{add\_rgb = False}.
  \item Enable the EMA normalization path in
        \texttt{\_process\_inferred\_depth} for GT depth.
  \item Concretely: force \texttt{\_use\_inferred\_depth = True} but
        provide simulator GT depth (raw meters) as the ``inferred'' source.
        This reuses the existing EMA pipeline on GT data.
  \item Alternative implementation: add a
        \texttt{force\_ema\_normalization} flag that applies EMA to the
        simulator's raw depth output \emph{before} its linear normalization.
\end{enumerate}

\subsection{\texttt{O2-gt-affine}: GT Depth with Random Affine Perturbation}

After the simulator's linear normalization, apply a per-batch random
affine:
\begin{align*}
  d' &= (1 + \epsilon_s) \cdot d + \epsilon_b \\
  \epsilon_s &\sim \mathcal{U}(-0.3, 0.3), \quad
  \epsilon_b \sim \mathcal{U}(-0.2, 0.2)
\end{align*}
Clamp result to $[-0.5, 0.5]$.  This tests CNN robustness to distribution shift.

\subsection{\texttt{N1-affine-fit}: Pre-computed Affine Normalization}

Replace EMA normalization with a fixed affine derived from least-squares
calibration:
\[
  d_{\text{GT}} \approx 0.088 \cdot d_{\text{inferred}} + 0.858
\]
Then apply standard GT normalization:
$d_{\text{norm}} = \text{clamp}(d_{\text{aligned}}, 0, 2) / 2 - 0.5$.

\subsection{\texttt{N5-running-meanstd}: Mean/Std Normalization}

Replace percentile-based EMA with running mean and standard deviation:
\[
  d_{\text{norm}} = \tanh\!\left(\frac{d - \mu_{\text{EMA}}}
  {2\,\sigma_{\text{EMA}}}\right) \cdot 0.5
\]
This preserves outlier structure (gaps produce values far from the mean)
while bounding the output to $[-0.5, 0.5]$.

\subsection{\texttt{D2-relative-inv}: Inverted Relative Model}

Use the relative DA-V2 model and invert the output:
$d' = \max(d_{\text{batch}}) - d_{\text{raw}}$, then apply EMA
normalization.  This corrects the inverted direction ($r = -0.59 \to
+0.59$).

\subsection{\texttt{A1-proprio-hist}: Proprioceptive History Encoder}

Add a 1D convolutional encoder over the last 10 proprioceptive frames
(following extreme-parkour \cite{cheng2023extreme}):
\begin{itemize}[nosep]
  \item Buffer: $(B, 10, 53)$ proprioceptive observations.
  \item 1D Conv: $53 \to 32$ channels, kernel 4, stride 2 $\to$ flatten
        $\to$ MLP $\to$ 32-dim embedding.
  \item Concatenate with depth latent before the actor MLP.
\end{itemize}

\subsection{\texttt{A3-2frame-input}: Two-Channel Depth Input}

Modify \texttt{DepthBackbone58x87} to accept 2-channel input:
\texttt{nn.Conv2d(2, 32, ...)} instead of \texttt{nn.Conv2d(1, 32, ...)}.
Pass both \texttt{\_depth\_buffer[:, 0]} (current) and
\texttt{\_depth\_buffer[:, 1]} (previous) as a stacked tensor.
This gives the CNN explicit access to temporal differences without
relying on the GRU.

\end{document}
